\subsection{Le module \emph{luadraw\_palettes}}

Le module \cmd{luadraw\_palettes}\footnote{Ce module est une contribution de \href{https://github.com/projetmbc/for-writing/tree/main/@prism}{Christphe BAL}.} définit $261$ palettes de couleurs portant chacune un nom. Une palette est une liste (table) de couleurs qui sont elles-mêmes des listes de trois valeurs numériques entre $0$ et $1$ (composantes rouge, verte et bleue). La liste de ces palettes ainsi que leur rendu, peuvent être visualisés dans ce \href{luadraw_palettes_list.pdf}{document}. 

\textbf{Ce module n'ajoute pas de nouvelles méthodes graphiques, mais il renvoie une table de définitions de couleurs}, plus la fonction \cmd{getPal(palname, options)} Par conséquent, il s'utilise ainsi (exemple) :
\begin{Luacode}
local pal = require 'luadraw_palettes'
local color = pal.Blackbody
local BlackbodyTransformed = pal.getPal( -- renvoie une nouvelle palette
    color, -- nom d'une palette
    {
    extract = {2, 5, 8, 9}, -- numéros des couleurs à extraire
    shift = 1, -- décalage parmi les couleurs extraites, ce qui donne ici: 5,8,9,2
    reverse = true -- inversion de l'ordre, ce qui donne ici: 2,9,8,5    
    }
)
\end{Luacode}
